\noindent\textbf{Por qué se usa este procedimiento.} Calcular primero $I_z$ de alma y alas. La ley $q=C_yQ/I_z$ se integra sobre el ala: $F=\int_0^b C_y(th/2)s/I_z\,ds=C_yhtb^2/(4I_z)$. No se multiplica el flujo máximo por toda la anchura.

\begin{enumerate}

\item La inercia es $t h^2 (h + 6b)/12$.

\end{enumerate}

\noindent\textbf{Operaciones y resultados intermedios.}

\begin{enumerate}

\item Fuerza en cada ala.

\[\fitmath{y_{1}=\frac{2000\cdot 100\cdot 2\cdot {50}^{2}}{4\cdot \frac{2\cdot {100}^{2}\cdot \left(100+6\cdot 50\right)}{12}}\approx 375\,\mathrm{N}}\]

\end{enumerate}
